\hypertarget{appendix-7a06e5baac04}{}\label{referencestart:appendix-7a06e5baac04}
\section{Copy-on-write}\label{appendix-7a06e5baac04-copy-on-write}

\hypertarget{appendix-7a06e5baac04-copy-on-write}{}

\textbf{Curriculum depth:} Introduced · \textbf{First seen:} C003 · \textbf{Latest development:} C003

\begin{allmightroledefinition}

\subsection{Definition and mechanism}\label{appendix-7a06e5baac04-definition-and-mechanism}

\hypertarget{appendix-7a06e5baac04-definition-and-mechanism}{}

COW initially shares physical contents while deferring private copies. A permitted write to a COW-protected mapping faults; OS allocates/copies a frame, updates the writer's mapping, and retries. Other mappings retain old contents. Not every read-only page is COW, so not every write fault is recoverable.

\end{allmightroledefinition}

\subsection{Prerequisites}\label{appendix-7a06e5baac04-prerequisites}

\hypertarget{appendix-7a06e5baac04-prerequisites}{}

\begin{itemize}
\tightlist
\item
  \hyperref[appendix-cc0909a7ef14]{Page faults and demand paging} --- Deferred private copying is implemented through a recoverable write fault.
\end{itemize}

\subsection{Used by}\label{appendix-7a06e5baac04-used-by}

\hypertarget{appendix-7a06e5baac04-used-by}{}

Update downstream uses when later chunks build on this concept.

\subsection{Recall and next depth}\label{appendix-7a06e5baac04-recall-and-next-depth}

\hypertarget{appendix-7a06e5baac04-recall-and-next-depth}{}

Show frame identity before and after one process writes.

\textbf{Pending:} Fork behavior, mapping flags, reference tracking and costs.

\subsection{Source and navigation}\label{appendix-7a06e5baac04-source-and-navigation}

\hypertarget{appendix-7a06e5baac04-source-and-navigation}{}

\hyperref[01-foundations-i-ch002-where-the-bytes-live]{C003 teaching chapter} · \hyperref[appendix-fd090d982a15-m003---parts-41-60]{Part definitions} · \href{../07-Vision-XR-and-Playable-Worlds-P01/main.pdf\#nameddest=appendix-4ad47bdd269a}{Dependency Map (Vision, XR, and Playable Worlds, Part 1, Reference notes: Dependency Map)}

\subsection{Reconstruction depth}\label{appendix-7a06e5baac04-reconstruction-depth}

\hypertarget{appendix-7a06e5baac04-reconstruction-depth}{}

Copy-on-write delays physical copying until shared data is actually modified. At the virtual-memory level, two address spaces may initially map the same physical page as read-only. A write triggers a protection fault; the operating system allocates a private page, copies old contents, updates the writer's page table and resumes the instruction.

The key invariant is logical independence after a write, even though physical storage was shared before it. This is why process creation can be cheap initially without eagerly duplicating every writable page.

COW has costs that disappear in a high-level description: the first write can incur a page fault and page copy, dirty-page accounting changes, and large working sets can turn an initially cheap clone into substantial later memory traffic. Huge pages and NUMA placement make those costs workload-dependent.

Do not generalize OS page COW into a recommendation for C++ value types. Library-level copy-on-write strings or containers complicate iterator/reference validity, threading and mutation semantics. The mechanism is useful when ownership and fault or copy behavior are explicit.
